% !TEX program = xelatex
\documentclass[10pt, a4paper]{article}

% ===== Page setup =====
\usepackage[left=1.0cm, right=1.0cm, top=0.6cm, bottom=0.6cm]{geometry}
\usepackage{titlesec}
\usepackage{enumitem}
\usepackage{tabularx}
\usepackage{xcolor}
\usepackage{hyperref}
\usepackage{fancyhdr}

% ===== Colors =====
\definecolor{primary}{RGB}{0, 90, 160}
\definecolor{darktext}{RGB}{40, 40, 40}
\definecolor{graytext}{RGB}{100, 100, 100}

% ===== Hyperlinks =====
\hypersetup{
    colorlinks=true,
    linkcolor=primary,
    urlcolor=primary,
    pdfborder={0 0 0}
}

% ===== No page numbers =====
\pagestyle{empty}

% ===== Paragraph settings =====
\setlength{\parindent}{0pt}
\setlength{\parskip}{0pt}

% ===== Section title style =====
\titleformat{\section}
    {\large\bfseries\color{primary}}
    {}
    {0em}
    {}
    [\vspace{-0.6em}\textcolor{primary}{\rule{\textwidth}{0.8pt}}]
\titlespacing*{\section}{0pt}{0.35em}{0.1em}

% ===== List style =====
\setlist[itemize]{
    leftmargin=1.5em,
    itemsep=0pt,
    parsep=0pt,
    topsep=0em,
    label=\textcolor{primary}{\textbullet}
}

% ===== Custom commands =====
% Entry: title | right-side information
\newcommand{\entry}[2]{%
    \vspace{0.1em}
    \noindent\textbf{#1}\hfill{\small\color{graytext}#2}\\[-0.4em]
}

% Subtitle row
\newcommand{\entrysub}[2]{%
    \noindent{\small #1}\hfill{\small\color{graytext}#2}\\[-0.5em]
}

\begin{document}

% ===== Header: name and contact information =====
\begin{center}
    {\LARGE\bfseries Yi Yuan}\\[0.3em]
    {\small
        \textbf{Email:} yuanyithu@126.com \quad
        \textbf{GitHub:} \href{https://github.com/yuanyithu}{github.com/yuanyithu} \\[0.2em]
        Male\enspace|\enspace Ph.D. Candidate in Physics, Tsinghua University
    }
\end{center}

% ===== Education =====
\section{Education}

\entry{Tsinghua University \enspace Ph.D. in Physics}{2023.08 -- 2028.06 (Expected)}
\entrysub{Department of Physics \enspace Ranked \textbf{1st} in Modern Physics in the Ph.D. Qualifying Exam \enspace Principles and Architecture of Quantum Computers (A+)}{Beijing}

\entry{B.S. in Physics, Fundamental Science Experimental Class (Mathematics and Physics), Tsinghua University}{2019.08 -- 2023.06}
\entrysub{Department of Physics \enspace GPA: 3.89/4.0 (Rank: 18/116) \enspace Outstanding Graduate \enspace Nominee, Ye Qisun Award}{Beijing}

% ===== Skills =====
\section{Technical Skills}
\begin{itemize}
    \item \textbf{Programming Languages:} Python (primary language for daily research and development)
    \item \textbf{AI/ML:} Familiar with Transformer architectures; experienced with JAX and PyTorch; hands-on experience in CNN modeling
    \item \textbf{Data and Computing:} NumPy vectorized computing, multiprocessing, data-driven modeling, and backtesting system development
    \item \textbf{Tools:} Git, Linux servers \quad \textbf{Language:} English -- proficient in reading and writing scientific literature
\end{itemize}

% ===== Publications =====
\section{Publications}
\noindent{\small \textbf{Yi Yuan}, Yuanchen Zhao, and Dong E. Liu. ``Zeno-Enhanced Probabilistic Error Cancellation with Quantum Error Detection Codes.'' \textit{arXiv preprint arXiv:2605.12149}, 2026. \href{https://arxiv.org/abs/2605.12149}{arXiv}}\\[-0.2em]
\noindent{\small Developed a perturbative QEDC+PEC protocol for near-term noisy quantum devices, using post-selection to reshape the effective noise channel and PEC to cancel residual logical errors.}

% ===== Research and projects =====
\section{Research and Projects}

\entry{Research in Quantum Computing and Quantum Information}{2022.09 -- Present}
\entrysub{Department of Physics, Tsinghua University}{Python, JAX, Tensor Network}
\begin{itemize}
    \item \textbf{Zeno-enhanced QEDC+PEC protocol} (2024 -- 2026, arXiv preprint): Combined quantum error-detecting codes with probabilistic error cancellation for near-term noisy quantum devices; in 200-qubit Iceberg-code logical GHZ preparation, reduced sampling overhead by \textbf{3--4 orders of magnitude} relative to standard PEC while maintaining $F\simeq0.956$
    \item \textbf{Classical Shadow error mitigation} (2024): Used tensor network methods to compute the error-mitigation complexity of classical shadows on noisy quantum circuits with $\log(n)$ depth, and analyzed its scaling with system size
    \item \textbf{Undergraduate thesis: complexity of quantum annealing} (2022.09 -- 2023.06): Based on the imaginary-time spacetime instanton action, explained why quantum annealing and classical QMC have matching complexity in ground-state energy search for double-well potentials, and identified the possibility of multi-path quantum speedup
\end{itemize}

\entry{Selected Projects}{}
\begin{itemize}
    \item \textbf{SEVI particle scattering analysis}: Built a PyTorch-based CNN to fit scattering-amplitude parameters, reaching the accuracy of analytical methods without injecting prior knowledge
    \item \textbf{LLM paper summarization tool} (2025): Automatically fetched daily arXiv papers, generated summaries through the Kimi API, and completed local deployment
\end{itemize}

% ===== Honors and leadership =====
\section{Honors and Leadership}
\noindent{\small Honors: 2021 Tsinghua Alumni -- Zhang Mingwei Scholarship \quad Multiple Department of Physics scholarships for academic excellence, innovation, and overall achievement in 2020--2021}\\[-0.2em]
\noindent{\small Leadership: 2023.08--2025.01 Student Counselor, Department of Physics, Tsinghua University \quad 2024--2025 Outstanding Student Leader}

\end{document}
